\item La velocidad de una ola de agua en aguas poco profundas se describe por $v = \sqrt{gd}$, donde $d$ es la profundidad del agua. Ya que cambia su velocidad, las olas se refractan cuando se mueven a una región de profundidad diferente.
\begin{enumerate}
    \item Trace conceptualmente la playa de un océano en el lado oriente de la masa de tierra. Muestre las líneas de contorno de profundidad constante bajo el agua si se trata de una pendiente razonablemente uniforme.
    \item Suponga que las olas se aproximan a la costa desde una tormenta situada a lo lejos hacia el norte-noreste. Demuestre que las olas se moverán en forma casi perpendicular a la orilla cuando lleguen a la playa.
    \item Trace un mapa de una costa con bahías y puntas (cabos) alternadas, como se sugiere en la figura P34.26, mostrando la forma de las líneas de contorno de profundidad constante.
    \item Suponga olas que se aproximan a la costa, que llevan energía con densidad uniforme a lo largo de frentes de onda originalmente rectos. Demuestre que la energía que llega a la orilla se concentra en las puntas y tiene menos intensidad en las bahías.
\end{enumerate}